\IfFileExists{stacks-project.cls}{%
\documentclass{stacks-project}
}{%
\documentclass{amsart}
}

% For dealing with references we use the comment environment
\usepackage{verbatim}
\newenvironment{reference}{\comment}{\endcomment}
%\newenvironment{reference}{}{}
\newenvironment{slogan}{\comment}{\endcomment}
\newenvironment{history}{\comment}{\endcomment}

% For commutative diagrams we use Xy-pic
\usepackage[all]{xy}

% We use 2cell for 2-commutative diagrams.
\xyoption{2cell}
\UseAllTwocells

% We use multicol for the list of chapters between chapters
\usepackage{multicol}

% This is generally recommended for better output
\usepackage{lmodern}
\usepackage[T1]{fontenc}

% For cross-file-references
\usepackage{xr-hyper}

% Package for hypertext links:
\usepackage{hyperref}

% For any local file, say "hello.tex" you want to link to please
% use \externaldocument[hello-]{hello}
\externaldocument[introduction-]{introduction}
\externaldocument[conventions-]{conventions}
\externaldocument[sets-]{sets}
\externaldocument[categories-]{categories}
\externaldocument[topology-]{topology}
\externaldocument[sheaves-]{sheaves}
\externaldocument[sites-]{sites}
\externaldocument[stacks-]{stacks}
\externaldocument[fields-]{fields}
\externaldocument[algebra-]{algebra}
\externaldocument[brauer-]{brauer}
\externaldocument[homology-]{homology}
\externaldocument[derived-]{derived}
\externaldocument[simplicial-]{simplicial}
\externaldocument[more-algebra-]{more-algebra}
\externaldocument[smoothing-]{smoothing}
\externaldocument[modules-]{modules}
\externaldocument[sites-modules-]{sites-modules}
\externaldocument[injectives-]{injectives}
\externaldocument[cohomology-]{cohomology}
\externaldocument[sites-cohomology-]{sites-cohomology}
\externaldocument[dga-]{dga}
\externaldocument[dpa-]{dpa}
\externaldocument[sdga-]{sdga}
\externaldocument[hypercovering-]{hypercovering}
\externaldocument[schemes-]{schemes}
\externaldocument[constructions-]{constructions}
\externaldocument[properties-]{properties}
\externaldocument[morphisms-]{morphisms}
\externaldocument[coherent-]{coherent}
\externaldocument[divisors-]{divisors}
\externaldocument[limits-]{limits}
\externaldocument[varieties-]{varieties}
\externaldocument[topologies-]{topologies}
\externaldocument[descent-]{descent}
\externaldocument[perfect-]{perfect}
\externaldocument[more-morphisms-]{more-morphisms}
\externaldocument[flat-]{flat}
\externaldocument[groupoids-]{groupoids}
\externaldocument[more-groupoids-]{more-groupoids}
\externaldocument[etale-]{etale}
\externaldocument[chow-]{chow}
\externaldocument[intersection-]{intersection}
\externaldocument[pic-]{pic}
\externaldocument[weil-]{weil}
\externaldocument[adequate-]{adequate}
\externaldocument[dualizing-]{dualizing}
\externaldocument[duality-]{duality}
\externaldocument[discriminant-]{discriminant}
\externaldocument[derham-]{derham}
\externaldocument[local-cohomology-]{local-cohomology}
\externaldocument[algebraization-]{algebraization}
\externaldocument[curves-]{curves}
\externaldocument[resolve-]{resolve}
\externaldocument[models-]{models}
\externaldocument[functors-]{functors}
\externaldocument[equiv-]{equiv}
\externaldocument[pione-]{pione}
\externaldocument[etale-cohomology-]{etale-cohomology}
\externaldocument[proetale-]{proetale}
\externaldocument[relative-cycles-]{relative-cycles}
\externaldocument[more-etale-]{more-etale}
\externaldocument[trace-]{trace}
\externaldocument[crystalline-]{crystalline}
\externaldocument[spaces-]{spaces}
\externaldocument[spaces-properties-]{spaces-properties}
\externaldocument[spaces-morphisms-]{spaces-morphisms}
\externaldocument[decent-spaces-]{decent-spaces}
\externaldocument[spaces-cohomology-]{spaces-cohomology}
\externaldocument[spaces-limits-]{spaces-limits}
\externaldocument[spaces-divisors-]{spaces-divisors}
\externaldocument[spaces-over-fields-]{spaces-over-fields}
\externaldocument[spaces-topologies-]{spaces-topologies}
\externaldocument[spaces-descent-]{spaces-descent}
\externaldocument[spaces-perfect-]{spaces-perfect}
\externaldocument[spaces-more-morphisms-]{spaces-more-morphisms}
\externaldocument[spaces-flat-]{spaces-flat}
\externaldocument[spaces-groupoids-]{spaces-groupoids}
\externaldocument[spaces-more-groupoids-]{spaces-more-groupoids}
\externaldocument[bootstrap-]{bootstrap}
\externaldocument[spaces-pushouts-]{spaces-pushouts}
\externaldocument[spaces-chow-]{spaces-chow}
\externaldocument[groupoids-quotients-]{groupoids-quotients}
\externaldocument[spaces-more-cohomology-]{spaces-more-cohomology}
\externaldocument[spaces-simplicial-]{spaces-simplicial}
\externaldocument[spaces-duality-]{spaces-duality}
\externaldocument[formal-spaces-]{formal-spaces}
\externaldocument[restricted-]{restricted}
\externaldocument[spaces-resolve-]{spaces-resolve}
\externaldocument[formal-defos-]{formal-defos}
\externaldocument[defos-]{defos}
\externaldocument[cotangent-]{cotangent}
\externaldocument[examples-defos-]{examples-defos}
\externaldocument[algebraic-]{algebraic}
\externaldocument[examples-stacks-]{examples-stacks}
\externaldocument[stacks-sheaves-]{stacks-sheaves}
\externaldocument[criteria-]{criteria}
\externaldocument[artin-]{artin}
\externaldocument[quot-]{quot}
\externaldocument[stacks-properties-]{stacks-properties}
\externaldocument[stacks-morphisms-]{stacks-morphisms}
\externaldocument[stacks-limits-]{stacks-limits}
\externaldocument[stacks-cohomology-]{stacks-cohomology}
\externaldocument[stacks-perfect-]{stacks-perfect}
\externaldocument[stacks-introduction-]{stacks-introduction}
\externaldocument[stacks-more-morphisms-]{stacks-more-morphisms}
\externaldocument[stacks-geometry-]{stacks-geometry}
\externaldocument[moduli-]{moduli}
\externaldocument[moduli-curves-]{moduli-curves}
\externaldocument[examples-]{examples}
\externaldocument[exercises-]{exercises}
\externaldocument[guide-]{guide}
\externaldocument[desirables-]{desirables}
\externaldocument[coding-]{coding}
\externaldocument[obsolete-]{obsolete}
\externaldocument[fdl-]{fdl}
\externaldocument[index-]{index}

% Theorem environments.
%
\theoremstyle{plain}
\newtheorem{theorem}[subsection]{Theorem}
\newtheorem{proposition}[subsection]{Proposition}
\newtheorem{lemma}[subsection]{Lemma}

\theoremstyle{definition}
\newtheorem{definition}[subsection]{Definition}
\newtheorem{example}[subsection]{Example}
\newtheorem{exercise}[subsection]{Exercise}
\newtheorem{situation}[subsection]{Situation}

\theoremstyle{remark}
\newtheorem{remark}[subsection]{Remark}
\newtheorem{remarks}[subsection]{Remarks}

\numberwithin{equation}{subsection}

% Macros
%
\def\lim{\mathop{\mathrm{lim}}\nolimits}
\def\colim{\mathop{\mathrm{colim}}\nolimits}
\def\Spec{\mathop{\mathrm{Spec}}}
\def\Hom{\mathop{\mathrm{Hom}}\nolimits}
\def\Ext{\mathop{\mathrm{Ext}}\nolimits}
\def\SheafHom{\mathop{\mathcal{H}\!\mathit{om}}\nolimits}
\def\SheafExt{\mathop{\mathcal{E}\!\mathit{xt}}\nolimits}
\def\Sch{\mathit{Sch}}
\def\Mor{\mathop{\mathrm{Mor}}\nolimits}
\def\Ob{\mathop{\mathrm{Ob}}\nolimits}
\def\Sh{\mathop{\mathit{Sh}}\nolimits}
\def\NL{\mathop{N\!L}\nolimits}
\def\CH{\mathop{\mathrm{CH}}\nolimits}
\def\proetale{{pro\text{-}\acute{e}tale}}
\def\etale{{\acute{e}tale}}
\def\QCoh{\mathit{QCoh}}
\def\Ker{\mathop{\mathrm{Ker}}}
\def\Im{\mathop{\mathrm{Im}}}
\def\Coker{\mathop{\mathrm{Coker}}}
\def\Coim{\mathop{\mathrm{Coim}}}

% Boxtimes
%
\DeclareMathSymbol{\boxtimes}{\mathbin}{AMSa}{"02}

%
% Macros for moduli stacks/spaces
%
\def\QCohstack{\mathcal{QC}\!\mathit{oh}}
\def\Cohstack{\mathcal{C}\!\mathit{oh}}
\def\Spacesstack{\mathcal{S}\!\mathit{paces}}
\def\Quotfunctor{\mathrm{Quot}}
\def\Hilbfunctor{\mathrm{Hilb}}
\def\Curvesstack{\mathcal{C}\!\mathit{urves}}
\def\Polarizedstack{\mathcal{P}\!\mathit{olarized}}
\def\Complexesstack{\mathcal{C}\!\mathit{omplexes}}
% \Pic is the operator that assigns to X its picard group, usage \Pic(X)
% \Picardstack_{X/B} denotes the Picard stack of X over B
% \Picardfunctor_{X/B} denotes the Picard functor of X over B
\def\Pic{\mathop{\mathrm{Pic}}\nolimits}
\def\Picardstack{\mathcal{P}\!\mathit{ic}}
\def\Picardfunctor{\mathrm{Pic}}
\def\Deformationcategory{\mathcal{D}\!\mathit{ef}}
\usepackage{cleveref}
\crefname{originalreference}{original reference}{original references}
\Crefname{originalreference}{Original reference}{Original references}
\makeatletter
\newlabel{algebra-lemma-NAK}{{00DV}{source}{}{sourcepointer-1}{}}
\newlabel{algebra-lemma-NAK@cref}{{[originalreference][0][]00DV}{[1][0][]source}{}{}{}}
\makeatother
\crefname{originalreference}{original reference}{original references}
\Crefname{originalreference}{Original reference}{Original references}
\makeatletter
\newlabel{algebra-lemma-flat-tor-zero}{{00HL}{source}{}{sourcepointer-2}{}}
\newlabel{algebra-lemma-flat-tor-zero@cref}{{[originalreference][0][]00HL}{[1][0][]source}{}{}{}}
\makeatother
\crefname{originalreference}{original reference}{original references}
\Crefname{originalreference}{Original reference}{Original references}
\makeatletter
\newlabel{algebra-lemma-smooth-at-point}{{07BU}{source}{}{sourcepointer-3}{}}
\newlabel{algebra-lemma-smooth-at-point@cref}{{[originalreference][0][]07BU}{[1][0][]source}{}{}{}}
\makeatother
\crefname{originallemma}{original lemma}{original lemmas}
\Crefname{originallemma}{Original lemma}{Original lemmas}
\makeatletter
\newlabel{lemma-injective}{{09PV}{source}{}{sourcepointer-4}{}}
\newlabel{lemma-injective@cref}{{[originallemma][0][]09PV}{[1][0][]source}{}{}{}}
\makeatother
\crefname{originalreference}{original reference}{original references}
\Crefname{originalreference}{Original reference}{Original references}
\makeatletter
\newlabel{more-algebra-definition-pseudo-coherent-perfect}{{067H}{source}{}{sourcepointer-5}{}}
\newlabel{more-algebra-definition-pseudo-coherent-perfect@cref}{{[originalreference][0][]067H}{[1][0][]source}{}{}{}}
\makeatother
\crefname{originalreference}{original reference}{original references}
\Crefname{originalreference}{Original reference}{Original references}
\makeatletter
\newlabel{more-algebra-lemma-Noetherian-pseudo-coherent}{{066E}{source}{}{sourcepointer-6}{}}
\newlabel{more-algebra-lemma-Noetherian-pseudo-coherent@cref}{{[originalreference][0][]066E}{[1][0][]source}{}{}{}}
\makeatother
\title{Perfect diagonal forces smoothness at the local level}
\author{The Stacks Project Authors}
\date{stacks-project:dpa.tex:local-snapshot-20261012; GNU FDL 1.2 or later}
\begin{document}
\maketitle
\noindent\textit{Editorial note.} Two standard identifications are left unverified in-span: the diagonal expansion ('a general fact about elements of A[x\_1, ..., x\_n] whose proof we omit') and 'Another computation shows that a\_ij is the image of the partial derivative' -- both are mechanical instantiations of the Jacobian presentation displayed two lines later; flagged for the reviewer.\par\medskip
\noindent\textit{Allowed standard prerequisites (editorial).}\begin{itemize}
\item Injectivity on the conormal level. lemma-injective (same chapter, named). Supplies I/mI -> J/mJ injective for the presentation data I subset J.
\item Smoothness from the naive cotangent complex. algebra-lemma-smooth-at-point (cross-chapter, named). Turns the freeness of NL\_\{C/A\} (x)\_A B into an open neighbourhood smooth over A.
\end{itemize}\medskip
\ifdefined\hypertarget\hypertarget{sourcepointer-1}{}\fi
\noindent\textit{Source reference.} Original reference 00DV: focused-seeds/geometry/S9-dpa-local-perfect-diagonal/ref-tags.txt, source lines 5--5. Navigation pointer only: the counted proof cites this statement by \textbackslash{}ref as a standing input; the deduction the counted proof has to make is carried out inside the counted span, and no part of that deduction is delegated to this pointer.\par\medskip
\ifdefined\hypertarget\hypertarget{sourcepointer-2}{}\fi
\noindent\textit{Source reference.} Original reference 00HL: focused-seeds/geometry/S9-dpa-local-perfect-diagonal/ref-tags.txt, source lines 6--6. Navigation pointer only: the counted proof cites this statement by \textbackslash{}ref as a standing input; the deduction the counted proof has to make is carried out inside the counted span, and no part of that deduction is delegated to this pointer.\par\medskip
\ifdefined\hypertarget\hypertarget{sourcepointer-3}{}\fi
\noindent\textit{Source reference.} Original reference 07BU: focused-seeds/geometry/S9-dpa-local-perfect-diagonal/ref-tags.txt, source lines 7--7. Navigation pointer only: the counted proof cites this statement by \textbackslash{}ref as a standing input; the deduction the counted proof has to make is carried out inside the counted span, and no part of that deduction is delegated to this pointer.\par\medskip
\ifdefined\hypertarget\hypertarget{sourcepointer-4}{}\fi
\noindent\textit{Source reference.} Original lemma 09PV: sources/stacks-dpa/pinned-original-source.tex, source lines 1790--1790. Navigation pointer only: the counted proof cites this statement by \textbackslash{}ref as a standing input; the deduction the counted proof has to make is carried out inside the counted span, and no part of that deduction is delegated to this pointer.\par\medskip
\ifdefined\hypertarget\hypertarget{sourcepointer-5}{}\fi
\noindent\textit{Source reference.} Original reference 067H: focused-seeds/geometry/S9-dpa-local-perfect-diagonal/ref-tags.txt, source lines 8--8. Navigation pointer only: the counted proof cites this statement by \textbackslash{}ref as a standing input; the deduction the counted proof has to make is carried out inside the counted span, and no part of that deduction is delegated to this pointer.\par\medskip
\ifdefined\hypertarget\hypertarget{sourcepointer-6}{}\fi
\noindent\textit{Source reference.} Original reference 066E: focused-seeds/geometry/S9-dpa-local-perfect-diagonal/ref-tags.txt, source lines 9--9. Navigation pointer only: the counted proof cites this statement by \textbackslash{}ref as a standing input; the deduction the counted proof has to make is carried out inside the counted span, and no part of that deduction is delegated to this pointer.\par\medskip
\begin{lemma}
\label{lemma-local-perfect-diagonal}
Let $A \to B$ be a local ring homomorphism of Noetherian local rings such that
$B$ is flat and essentially of finite type over $A$. If
$$
B \otimes_A B \longrightarrow B
$$
is a perfect ring map, i.e., if $B$ has finite tor dimension over
$B \otimes_A B$, then $B$ is the localization of a smooth $A$-algebra.
\end{lemma}

\begin{proof}
As $B$ is essentially of finite type over $A$, so is $B \otimes_A B$ and
in particular $B \otimes_A B$ is Noetherian. Hence the quotient $B$ of
$B \otimes_A B$ is pseudo-coherent over $B \otimes_A B$
(More on Algebra, Lemma \ref{more-algebra-lemma-Noetherian-pseudo-coherent})
which explains why perfectness of the ring map (More on Algebra, Definition
\ref{more-algebra-definition-pseudo-coherent-perfect}) agrees with the
condition of finite tor dimension.

\medskip\noindent
We may write $B = R/K$ where $R$ is the localization of $A[x_1, \ldots, x_n]$
at a prime ideal and $K \subset R$ is an ideal. Denote
$\mathfrak m \subset R \otimes_A R$ the maximal ideal which is the inverse
image of the maximal ideal of $B$ via the surjection
$R \otimes_A R \to B \otimes_A B \to B$. Then we have surjections
$$
(R \otimes_A R)_\mathfrak m \to (B \otimes_A B)_\mathfrak m \to B
$$
and hence ideals $I \subset J \subset (R \otimes_A R)_\mathfrak m$
as in Lemma \ref{lemma-injective}. We conclude that
$I/\mathfrak m I \to J/\mathfrak m J$ is injective.

\medskip\noindent
Let $K = (f_1, \ldots, f_r)$ with $r$ minimal. We may and do assume that
$f_i \in R$ is the image of an element of $A[x_1, \ldots, x_n]$ which we
also denote $f_i$. Observe that $I$ is generated
by $f_1 \otimes 1, \ldots, f_r \otimes 1$ and
$1 \otimes f_1, \ldots, 1 \otimes f_r$. We claim that this is a minimal
set of generators of $I$. Namely, if $\kappa$ is the common residue field
of $R$, $B$, $(R \otimes_A R)_\mathfrak m$, and $(B \otimes_A B)_\mathfrak m$
then we have a map
$R \otimes_A R \to R \otimes_A \kappa \oplus \kappa \otimes_A R$
which factors through $(R \otimes_A R)_\mathfrak m$. Since $B$ is
flat over $A$ and since we have the short exact sequence
$0 \to K \to R \to B \to 0$ we see that
$K \otimes_A \kappa \subset R \otimes_A \kappa$, see
Algebra, Lemma \ref{algebra-lemma-flat-tor-zero}.
Thus restricting the map
$(R \otimes_A R)_\mathfrak m \to R \otimes_A \kappa \oplus \kappa \otimes_A R$
to $I$ we obtain a map
$$
I \to K \otimes_A \kappa \oplus \kappa \otimes_A K \to
K \otimes_B \kappa \oplus \kappa \otimes_B K.
$$
The elements
$f_1 \otimes 1, \ldots, f_r \otimes 1, 1 \otimes f_1, \ldots, 1 \otimes f_r$
map to a basis of the target of this map, since by Nakayama's lemma
(Algebra, Lemma \ref{algebra-lemma-NAK})
$f_1, \ldots, f_r$ map to a basis of $K \otimes_B \kappa$.
This proves our claim.

\medskip\noindent
The ideal $J$ is generated by $f_1 \otimes 1, \ldots, f_r \otimes 1$
and the elements $x_1 \otimes 1 - 1 \otimes x_1, \ldots,
x_n \otimes 1 - 1 \otimes x_n$ (for the proof it suffices to
see that these elements are contained in the ideal $J$).
Now we can write
$$
f_i \otimes 1 - 1 \otimes f_i =
\sum g_{ij} (x_j \otimes 1 - 1 \otimes x_j)
$$
for some $g_{ij}$ in $(R \otimes_A R)_\mathfrak m$. This is a general
fact about elements of $A[x_1, \ldots, x_n]$ whose proof we omit.
Denote $a_{ij} \in \kappa$ the image of $g_{ij}$. Another computation
shows that $a_{ij}$ is the image of $\partial f_i / \partial x_j$ in $\kappa$.
The injectivity of $I/\mathfrak m I \to J/\mathfrak m J$ and the remarks
made above forces the matrix $(a_{ij})$ to have maximal rank $r$.
Set
$$
C = A[x_1, \ldots, x_n]/(f_1, \ldots, f_r)
$$
and consider the naive cotangent complex
$\NL_{C/A} \cong (C^{\oplus r} \to C^{\oplus n})$
where the map is given by the matrix of partial derivatives.
Thus $\NL_{C/A} \otimes_A B$
is isomorphic to a free $B$-module of rank $n - r$ placed in degree $0$.
Hence $C_g$ is smooth over $A$ for some $g \in C$ mapping to a unit
in $B$, see Algebra, Lemma \ref{algebra-lemma-smooth-at-point}.
This finishes the proof.
\end{proof}

\end{document}
